\section{Introduction}
\label{sec:introduction}

Software-engineering agents increasingly operate across repositories, shells, test runners, browsers, issue trackers, build systems, and other development tools. SWE-bench established repository-level issue resolution as a measurable software-engineering task \citep{jimenez2024swebench}; SWE-agent showed that the design of the agent--computer interface materially affects a language model's ability to navigate repositories, edit files, and execute tests \citep{yang2024sweagent}; and OpenHands provides a general platform and later SDK for software agents with sandboxed execution, lifecycle control, memory management, and model routing \citep{wang2024openhands,wang2025openhandsSDK}. These systems demonstrate substantial progress toward agents that act like software developers. They do not settle a different systems question: \emph{what must remain persistent when engineering work extends beyond one reasoning session?}

The conventional framing makes the intelligent developer the persistent object. A long-running agent accumulates conversation, observations, summaries, working memory, execution state, and tool history so that later decisions can depend on earlier reasoning. This is useful and sometimes necessary. Modern agent-memory research explicitly studies how experience should be stored, consolidated, retrieved, reflected upon, and reused \citep{du2026memory,luo2026memory}. Context-engineering practice likewise treats the model-visible token budget as a resource that must be curated over long interactions \citep{anthropic2025context}. RaS does not argue that these mechanisms are universally unnecessary.

It asks whether software engineering permits a stronger inversion:

\begin{quote}
\textbf{The persistent object need not be the developer; it can be the software.}
\end{quote}

Software projects already possess unusually rich domain-native durable state. Source code captures implementation. Tests encode executable behavioural constraints. Schemas and interfaces capture contracts. Configuration and build definitions capture environmental assumptions. Architecture records and documentation can preserve rationale and non-obvious constraints. Git history captures accepted transitions. Continuous-integration evidence can attest that particular states compiled, tested, or satisfied external checks. A software project therefore has a natural substrate into which the consequences of prior reasoning can be externalised.

\emph{Repository-as-State} (RaS) investigates whether that substrate can carry enough engineering continuity that the high-capability reasoning process no longer needs to remain alive between bounded tasks. The central proposition is:

\begin{quote}
\textbf{Durable engineering continuity does not necessarily need to live inside the expensive high-capability reasoning process.}
\end{quote}

The distinction between “cheap” and “disposable” is central. RaS does not primarily attempt to make high-quality reasoning cheap. It attempts to make high-quality reasoning \emph{disposable}. A high-capability reasoning process may reconstruct relevant engineering state, reason deeply, produce a bounded transition or work package, delegate routine execution, interpret deterministic evidence, persist the accepted result into the repository, and then cease to exist. A later completely fresh reasoner should need the updated repository and the new task, rather than the full conversational history of its predecessors.

Figure~\ref{fig:ras-architecture} summarises the proposed loop.

\begin{figure}[htbp]
\centering
\input{figures/ras-architecture}
\caption{Conceptual Repository-as-State architecture. The durable object is repository state; reasoning and execution processes may be replaced between validated transitions. No empirical performance is implied.}
\label{fig:ras-architecture}
\end{figure}

The architecture can be written at a high level as
\[
R_t + U_t
\rightarrow
\text{bounded reconstruction}
\rightarrow
\text{high-capability reasoning}
\rightarrow
\text{execution}
\rightarrow
\text{validation}
\rightarrow
R_{t+1}.
\]
After the state transition, the reasoning process may be destroyed. This turns continuity into a property of the engineering programme rather than a property of one long-lived model session.

The idea arose from repeated practical software-engineering workflows in which difficult reasoning was performed in high-capability sessions, detailed implementation packages were handed to comparatively shallow execution agents, and compilers, tests, runtime observations, documentation, and Git persisted accepted results. Fresh reasoning sessions could often resume from repository state. That observation is valuable motivation because the architecture was recognised through use rather than invented purely as an abstract construction. It is nevertheless \emph{observed motivation}, not controlled empirical proof.

RaS therefore frames seven contributions at paper v0.1:

\begin{enumerate}
    \item a Repository-as-State architecture in which durable engineering state is primary and high-capability reasoning is potentially ephemeral;
    \item a theoretical formulation of approximate repository sufficiency;
    \item an explicit separation between durable repository state and the bounded reasoning view reconstructed for a task;
    \item a cumulative-context and provider-neutral cost model under stated assumptions, with reconstruction cost treated as a central falsifier;
    \item Forced-State-Reset Evaluation as a methodology for deliberately removing conversational continuity;
    \item the proposed Repository Resumability Index (RRI) and a separate state-reconstruction probe for analysing resumption;
    \item a secondary tiered-execution and least-privilege architecture that separates reasoning capability from routine operational permissions.
\end{enumerate}

The paper does not claim empirical findings. P0, the initial Forced-State-Reset Pilot, has not been executed. The objective of v0.1 is to define the theory and evaluation clearly enough that subsequent evidence can strengthen, weaken, narrow, or falsify the hypotheses without retroactively changing what was meant by RaS.
